% Generated draft: requires mathematical and source review.
\documentclass[11pt]{article}
\usepackage[margin=25mm]{geometry}
\usepackage{fontspec}
\setmainfont{Latin Modern Roman}
\usepackage{amsmath,amssymb,mathtools,unicode-math}
\setmathfont{Latin Modern Math}
\usepackage{enumitem,needspace,xurl,hyperref}
\setlength{\parindent}{0pt}
\setlength{\parskip}{5pt}
\setlength{\emergencystretch}{3em}
\newcommand{\N}{\mathbb{N}}
\newcommand{\Z}{\mathbb{Z}}
\newcommand{\Q}{\mathbb{Q}}
\newcommand{\R}{\mathbb{R}}
\newcommand{\C}{\mathbb{C}}
\newcommand{\F}{\mathbb{F}}
\newcommand{\A}{\mathbb{A}}
\newcommand{\PP}{\mathbb{P}}
\newcommand{\E}{\mathbb{E}}
\newcommand{\eps}{\varepsilon}
\newcommand{\dd}{\,\mathrm{d}}
\begin{document}
\section*{UnsolvedMath ID: 7600018}
\textbf{Source problem code:} catalogue entry
\subsection{Definitions and mathematical statement}
Let \(T\) be a complete consistent first-order theory in a language \(L\). A model \(M\models T\) is atomic if, for every finite tuple \(\bar a\in M\), its complete type over the empty set is isolated: there is \(\varphi(\bar x)\in\operatorname{tp}(\bar a/\varnothing)\) such that
\[
T\vdash\forall\bar x\,(\varphi(\bar x)\to\psi(\bar x))
\quad\text{for every }\psi\in\operatorname{tp}(\bar a/\varnothing).
\]
The atomic class is categorical in an infinite cardinal \(\kappa\) if it has a model of size \(\kappa\) and any two atomic models of that size are isomorphic. Existence is part of categoricity, so an empty size class is not counted vacuously.

On the all-finite-index reading of the phrase the \(\aleph_n\), the proposed transfer is
\[
\left[\forall n<\omega,\quad\operatorname{At}(T)
\text{ is categorical in }\aleph_n\right]
\quad\Longrightarrow\quad
\left[\forall\kappa\text{ infinite},\quad\operatorname{At}(T)
\text{ is categorical in }\kappa\right].
\]
The source does not explicitly state whether it means all finite \(n\), one fixed \(n\), or only positive indices. Nor does it specify a countable language or additional set-theoretic assumptions. Those choices matter and are not silently inserted. Every cardinal is interpreted in the standard categoricity-transfer sense of every infinite cardinal; literal inclusion of finite cardinals would require a different conclusion for theories already having infinite models.

This is about the restricted class of atomic models, not categoricity of all models of \(T\), and not an automatic application of the ordinary first-order categoricity theorem. The supplied research summary reports partial transfers under extra hypotheses, including excellence/tameness-like conditions, and no verified full general result. Such additional hypotheses are not added to the question or presumed for every atomic class; references to minimal models and abstract elementary classes do not change the class being quantified.
\subsection{Short English statement}
If a theory's atomic models are categorical at all intended \(\aleph_n\), must they be categorical at every infinite cardinal? Fix the index quantifier, language size, and set-theoretic background before claiming a categoricity transfer.
\subsection{Sources}
\begin{itemize}
\item Local UnsolvedMath registry, permanent ID \(7600018\), exact source code \texttt{catalogue entry}, atomic-model categoricity question.
\item Wikipedia, \emph{List of unsolved problems in mathematics}, pinned revision 1366636767, Model theory and formal languages, bullet 18; \url{https://en.wikipedia.org/w/index.php?title=List_of_unsolved_problems_in_mathematics&oldid=1366636767#Model_theory_and_formal_languages}. MediaWiki-source extraction; recorded access 2026-07-29.
\item Supplied PARTIAL-PROGRESS research summary names Lessmann, Baldwin--Kolesnikov, and Shelah, including arXiv:0903.3428, for qualified atomic/excellent-class transfer results. It supplies no complete references or exact hypotheses matching the unrestricted question.
\end{itemize}
\end{document}
